% Transcription — fonds Alexandre Grothendieck, shelfmark 10, batch 4 (pages 61-80).
% Established from the facsimile archives/batches/10/batch-04.pdf (a mirror of
% https://grothendieck.umontpellier.fr/10.pdf), per the skill
% transcribe-grothendieck.
%
% Pass: Opus 5 (claude-opus-5), 2026-09-12 — first pass, unchecked against the
% pages by a human.
%
% Every one of the twenty pages is mathematics and every one is on the subject
% of the carton, which the three previous batches were almost entirely not.
% The lot holds two typescripts, neither of them his, both corrected in his
% hand, and they run without a break from the first page to the last.
%
% Pages 61-65, « Etude numérique des catégories tannakiennes », in three
% lettered parts:
%
%   61-62 a) the numerical invariants of a k-linear abelian category with
%         finite-dimensional Hom: for M simple the subcategory it generates is
%         that of finite-dimensional vector spaces over K = End(M), so M stays
%         semi-simple under k'/k exactly when K ⊗_k k' is a product of matrix
%         algebras, and absolutely semi-simple exactly when the centre Z of K
%         is separable over k; K(𝓜) free on the set Σ(𝓜) of classes of simple
%         objects; the decomposition of M_k' read off K ⊗_k k' = ∏ M_{n_i}(K'_i)
%         with Σ(𝓜) ≃ Σ(𝓜')/π for k'/k quasi-galois; and σ ↦ (Z_σ, ξ_σ ∈ Br(Z_σ)),
%         with the transition map K(𝓜) → K(𝓜_{k'}), σ ↦ n_i(σ,i), n_i = d/d_i
%         the ratio of the orders of ξ_σ and of its image.
%   62-63 b) the four « caractères numériques » of a ⊗-category — the set Σ, the
%         multiplication of K(𝓜) through the structure constants c^ρ_{σ,τ}, the
%         map σ ↦ (Z_σ, ξ_σ), and the rank function 𝓜 → Z — then K^eff and the
%         remark that Σ is recovered from it as the set of non-zero minimal
%         elements; the reduced characteristic polynomials det_M f, Tr_M f,
%         Pol_M(f,t) in the Azumaya algebras End(M_i); and the warning that the
%         numerical characters do not in general determine the category, even
%         over an algebraically closed k of characteristic 0.
%   64-65 c) the case of a band of multiplicative type: the numerical data
%         yields G = D(Γ) with Γ = Σ' carrying the action of π, u_σ: G → Z_σ^*
%         = ∏_{Z/k} G_m with u_σ(ξ) = ξ_σ, and n(σ) = card(σ)·ord(ξ_σ); then the
%         question whether ξ ∈ H²(k_pl, G) is known from the u_σ(ξ), settled
%         for G diagonalisable, and Exemple 1 (a one-dimensional torus, by
%         Hilbert 90). Exemple 2, on G' = Ker n·id_G, is cancelled by long
%         diagonals and the run breaks off with it.
%
% Pages 66-80, « Catégories tannakiennes définies par des cristaux », in nine
% numbered sections:
%
%   66-68 1-4: crystals on the absolute crystalline site of a field k of
%         characteristic p, described through a p-ring W as finite-type modules
%         with a flat absolute connection; isocrystals ≃ Modf(K); F-crystals
%         F_M: M^σ → M and non-degeneracy; FCrisö(k), the F-isocristaux
%         effectifs, the inverse Tate object K(-1), and FCriso(k) obtained by
%         inverting it; for k perfect the interpretation as (E, F_E) with F_E a
%         σ-linear automorphism, effective iff the iterates F_E^n are bounded;
%         and the gerbe G(k) with its effectivity structure and natural fibre
%         functor over K(k).
%   69-70 the crystalline cohomology functor H*_cris on proper smooth schemes,
%         additive, compatible with ⊗ and with the Koko rule for the
%         commutativity constraint, H²_cris(P¹) = W(-1) and H^{2d}_cris(X) ≃
%         W(-d); the gerbe homomorphism ∏(k)_{Q_p} ← G(k) induced by
%         Mot(k) → FCriso(k); ζ_X computed from H*(X) for k = F_q; then 5:
%         F-crystals of slope zero, M = V ⊗_{Z_p} W, Proposition 5.1 (V ↦ V ⊗ W
%         is an equivalence onto the F-crystals with F_M an isomorphism) and
%         Corollaire 5.1.
%   71-73 the four equivalent conditions for an (E,F_E) to be of slope zero,
%         the epimorphism of gerbes G(k) → ∏_p(k) and G(k) → ∏_p(k) × [G_m];
%         6: the behaviour under k → k', Proposition 6.1 for k'/k finite étale
%         with the exact sequence 0 → G'(k') → G(k) → [g] → 0, and Question 6.2
%         for k'/k galois possibly infinite; 7: the affirmative answer for k
%         finite, G(k) ≃ G(F_p) ×_{∏_p(F_p)} ∏_p(k), and the two fibre functors
%         over K whose comparison gives the extra structure of G(k).
%   74-76 8: the gerbe G(F_p) — its lien is the Q_p-pro-algebraic envelope of
%         the discrete group Z, hence G_u × G_r with Γ = Q̄_p^*; the quotient
%         G' by the subgroup of p-adic units, giving 0 → H → G → G' → 0 with
%         H = D(Q) diagonalisable through v_p, and the Tate augmentation
%         Z ↪ Q̄_p^*, Z ↪ Q; 8.2: the gerbe homomorphism i: G(F_p) → ∏_p(F_p),
%         the identity on liens, the torseur under G' that defines it, and the
%         canonical P_o attached to the maximal unramified extension of Q_p.
%   77-78 8.3: H the kernel gerbe, ξ ∈ H²(Q_p, H) = Hom(Q, Br(Q_p)) =
%         Hom(Q, Q/Z), ξ = ∂(η) with η ∈ H¹(Q_p, G') the class of P, the
%         quotient H' of G_r by the roots of unity and the roots of p, the
%         two-row diagram of extensions, and the check that ξ corresponds to
%         Q → Q/Z under Br(Q_p) ≃ Q/Z — « sauf erreur de signe, car je n'ai
%         nulle part fait attention aux signes »; 8.4: G(F̄_p) → G(F_p)
%         factors through H, and by Dieudonné-Manin it is an equivalence;
%         8.5: question 6.2 settled for k = F_p.
%   79-80 8.6: Rep(H), the slope of a simple object, the complete catalogue
%         E_{r,s} of rank r with End = K_{r,s} of class s/r in Br(Q_p) = Q/Z,
%         Rep(H) as the inductive limit of the Rep(H,r), and two exercises he
%         did not do — the ⊗-structure in terms of the E_{r,s}, and the
%         Dieudonné modules of the reduced Lubin-Tate groups; 9: for k
%         arbitrary the canonical FCriso(k) → Rep(H), H → G(k), the definition
%         of slope λ, the centrality of H → G(k) for k perfect and the
%         isopentic decomposition E = Σ E_λ, and E = Σ_{i∈Z} E_i(-i) with the
%         E_i of slope 0 ≤ i < 1.
%
% The machine had no script letters, no Greek and no bars: 𝓜, the group
% letters, the Π of the fundamental gerbe and the hats and underlinings that
% distinguish a scheme from a group are inked throughout, and page 71 carries
% a note of his to himself, « Π majuscule ? », about exactly that. Underlined
% capitals are his own convention, stated on page 64: « les soulignés
% indiquent qu'on prend des schémas sur k ». Insertions are given in \add{},
% deletions in \struck{}, his margin notes in \marginal{}; the machine's own
% overtyped corrections are not recorded one by one.
%
% Two readings are flagged rather than repaired. The hand's last line of page
% 62 writes the structure constants σ^ρ_{σ,τ} where the typed formula two
% lines above has c^ρ_{σ,τ}, and the note says so. On page 79 the margin's
% E_λ ⊗ E_λ' carries a subscript on the right-hand side that the leaf does not
% determine between λλ' and λ+λ'; it is given as written and marked.
%
% The dating is the inventory's own, for the whole shelfmark.
\documentclass[11pt,a4paper]{article}
% Le préambule commun à toutes les transcriptions du fonds Grothendieck.
%
% Il tient en une idée : **l'incertitude se compose, elle ne se résout pas.**
% Une transcription qui lisse un mot douteux a détruit la seule chose pour
% laquelle on la faisait — un lecteur doit pouvoir savoir, à chaque endroit,
% ce qui était sur le feuillet et ce qui a été ajouté. Les six macros qui
% suivent sont donc l'appareil critique, et non de la décoration.
%
% Les mêmes commandes sont reconnues par `scripts/render.mjs`, qui produit la
% vue de lecture du site. Une macro ajoutée ici doit l'être aussi là-bas,
% faute de quoi le rendu échouera bruyamment — ce qui est voulu : un rendu
% silencieusement faux est pire qu'un rendu absent.

\ProvidesPackage{grothendieck}[2026/08/08 Transcriptions du fonds Grothendieck]

\usepackage{amsmath,amssymb,amsthm}
\usepackage{mathrsfs}
\usepackage{stmaryrd}
% Les diagrammes commutatifs sont partout dans ce fonds : c'est la langue
% même de la géométrie algébrique de Grothendieck, et une transcription qui
% les rendrait en prose ne transcrirait rien.
\usepackage{tikz-cd}
% Français par défaut — c'est la langue des feuillets — mais l'anglais est
% chargé aussi : la traduction et le résumé sont des documents anglais, et la
% typographie française (espaces avant les deux-points) y serait une faute.
% Ces documents basculent par \selectlanguage{english} en tête de corps.
\usepackage[english,french]{babel}
\usepackage{fontspec}
\usepackage{geometry}
\geometry{a4paper,margin=25mm,marginparwidth=28mm}
\usepackage{marginnote}
\usepackage{xcolor}
% ulem, not soul. `soul`'s \st and \ul are built on a character-by-character
% scan that XeLaTeX's font machinery breaks: the document fails with an
% undefined control sequence rather than degrading. ulem does the same two jobs
% and is Unicode-engine safe. `normalem` stops it hijacking \emph, which the
% transcriptions use for Grothendieck's own emphasis.
\usepackage[normalem]{ulem}
\usepackage{hyperref}
\hypersetup{colorlinks=true,linkcolor=black,urlcolor=black,pdfborder={0 0 0}}

% --- Métadonnées du lot -----------------------------------------------------
% Reprises telles quelles par le script de rendu, qui les affiche en tête de
% la vue de lecture. Elles fixent ce que le fichier prétend couvrir : sans
% elles, une transcription flotte sans point d'attache dans un fonds de seize
% mille feuillets.

\newcommand{\folder}[1]{\gdef\grothFolder{#1}}
\newcommand{\batch}[1]{\gdef\grothBatch{#1}}
\newcommand{\leaves}[2]{\gdef\grothFirst{#1}\gdef\grothLast{#2}}
\newcommand{\foldertitle}[1]{\gdef\grothFoldertitle{#1}}
\gdef\grothFolder{?}\gdef\grothBatch{?}\gdef\grothFirst{?}\gdef\grothLast{?}\gdef\grothFoldertitle{}

% --- Le résumé --------------------------------------------------------------
%
% Il ouvre la lecture modernisée, sous le titre « Résumé », et il est composé
% en petit. Ce n'est pas un ornement : c'est le seul passage du document qui
% ne soit pas des mathématiques, et un lecteur doit pouvoir le reconnaître
% avant de l'avoir lu — puis le sauter s'il connaît déjà le sujet, ou s'y
% arrêter s'il ne le connaît pas. Le filet qui le suit marque la reprise du
% corps.

\newenvironment{resume}
  {\small\setlength{\parskip}{0.45em}}
  {\par\medskip\hrule\medskip}

% --- Mots-clés ---------------------------------------------------------------
%
% En anglais, en clôture du résumé de la lecture modernisée : le vocabulaire
% moderne sous lequel chercher ce que les feuillets construisent. Le manifeste
% du site (`npm run manifest`) les extrait pour étiqueter la cote entière —
% cette ligne est la source unique des tags, il n'y a pas de fichier à côté.
\newcommand{\keywords}[1]{\par\smallskip\noindent
  {\footnotesize\textsc{Keywords} --- \textit{#1}}\par}

% --- L'appareil critique ----------------------------------------------------

% Le numéro de feuillet, en marge. C'est l'ancre qui permet de régler un
% désaccord sur une formule en regardant *un* feuillet plutôt qu'une chemise
% de six cents. Le site s'en sert aussi pour tourner les pages du fac-similé
% quand on fait défiler la transcription.
\newcommand{\leaf}[1]{%
  \marginnote{\footnotesize\color{gray!70}\textsf{#1}}%
  \label{leaf:#1}\ignorespaces}

% Illisible : on ne devine pas. Un mot inventé qui se lit comme les autres est
% le pire résultat possible.
\newcommand{\ill}{\textcolor{red!60!black}{[\,\ldots\,]}}

% Lecture douteuse : le mot est donné, et signalé comme incertain.
\newcommand{\uncertain}[1]{\uline{#1}}

% Ajout de l'éditeur — une lettre manquante, un mot sous-entendu.
\newcommand{\add}[1]{\textcolor{blue!50!black}{[#1]}}

% Biffé par Grothendieck lui-même. Ce qu'il a rayé fait partie du document :
% une rature dit souvent où la pensée a changé de direction.
\newcommand{\struck}[1]{\sout{#1}}

% Note du transcripteur, sur l'état matériel du feuillet.
\newcommand{\note}[1]{\par\smallskip{\small\color{gray!60!black}\textit{[#1]}}\par\smallskip}

% Note marginale de Grothendieck, distincte du corps du texte.
\newcommand{\marginal}[1]{\par\smallskip\noindent\hspace{1em}%
  {\small\color{gray!60!black}#1}\par\smallskip}

% --- Titre ------------------------------------------------------------------

\AtBeginDocument{%
  \begin{center}
    {\large\bfseries Fonds Alexandre Grothendieck --- cote n\textsuperscript{o}~\grothFolder}\\[2pt]
    {\itshape\grothFoldertitle}\\[4pt]
    {\small feuillets \grothFirst--\grothLast\ (lot \grothBatch)}\\[2pt]
    {\footnotesize\color{gray!60!black}Transcription établie d'après le fac-similé numérisé
     par l'Université de Montpellier.\\
     Les lectures incertaines sont soulignées, les illisibles notées [\,\ldots\,],
     les ajouts entre crochets.}
  \end{center}
  \vspace{1em}}

\endinput


\folder{10}
\batch{4}
\pages{61}{80}
\dating{[à partir de 1958]}
\watermark{Édition de démonstration}
\foldertitle{Catégories tensorielles : notes manuscrites (s.d.), tapuscrits (s.d.)}

\begin{document}

\section*{Étude numérique des catégories tannakiennes (pages 61 à 65)}

\note{Tapuscrit, corrigé de sa main au crayon et à l'encre bleue. La lettre
$\mathcal{M}$, les lettres de groupes, les $\Sigma$, les $\sigma$ et les
$\xi$ manquaient à la machine et sont encrés partout ; on ne les relève pas
un à un. Les majuscules soulignées sont sa convention, énoncée page 64 :
« les soulignés indiquent qu'on prend des schémas sur $k$ ». Ses insertions
sont données entre crochets, ses suppressions barrées, ses notes marginales
en retrait.}

\page{61}
\emph{Etude numérique des catégories tannakiennes.}
\note{Sic, sans accent, comme sur la page.}

a) \quad Soit $\mathcal{M}$ une catégorie abélienne $k$-linéaire avec
$\operatorname{Hom}$ de dimension finie ($k$ un corps). Si $M$ est un objet
simple de $\mathcal{M}$, la sous-catégorie abélienne de $\mathcal{M}$ qu'il
engendre est équivalente à la catégorie des vectoriels de dimension finie sur
le corps $K = \operatorname{End}(M)$. On voit ainsi que l'objet $M$ reste
semi-simple par une extension $k'/k$ sss $K \otimes_{k} k'$ est un composé
\struck{\ill{}} d'algèbres de matrices (i.e. est une algèbre semi-simple).
Ceci est vrai \add{\ill{}} pour tout $k'$ sss c'est vrai pour
\struck{$k'$ clôture} \add{une ext.} \uncertain{parfaite} de $k$,
\struck{ou} sss le centre \struck{de} $Z$ de $K$ est séparable sur $k$. On
dit alors que $M$ est absolument semi-simple (à généraliser au cas
\add{où $M$ n'}est pas supposé simple).
\note{L'insertion au-dessus de « vrai » est une seule lettre, qui ne se lit
pas ; « parfaite » est souligné sur la page.}

Le groupe $K(\mathcal{M})$ est le groupe libre engendré par l'ensemble
$\Sigma(\mathcal{M})$ des classes d'objets simples de $\mathcal{M}$. Si $k'$
est une extension de $k$, la façon dont un objet simple $M$ de $\mathcal{M}$
tel que $M_{k'}$ soit semi-simple se décompose en objets simples se voit sur
la structure de $K \otimes_{k} k' = K'$ : si $K'$ est le produit d'algèbres
$M_{n_{i}}(K'_{i})$, où les $K'_{i}$ sont des corps gauches,
\add{$1 \leq i \leq r$,} alors $M_{k'}$ se décompose en $r$ composants
isotypiques (correspondant aux $K'_{i}$) chacun ayant $n_{i}$ composants
simples \add{$M'_{i}$}. Si $k'$ est une extension quasi-galoisienne de $k$,
les classes des $M'_{i}$ sont conjuguées entre elles par l'action de
$\operatorname{Gal}(k'/k) = \pi$. On trouve alors une bijection canonique
(si l'hypothèse faite est vérifiée pour tout objet simple \struck{de} $M$ de
$\mathcal{M}$, p.\,ex. si $k'/k$ est même galoisienne)
\[
\Sigma(\mathcal{M}) \simeq \Sigma(\mathcal{M}')/\pi .
\]

Il y aurait lieu de généraliser au cas où on ne suppose pas nécessairement
que les $K \otimes_{k} k'$ sont semi-simples …

\add{Si $\mathcal{M}$ est semi-simple,} La catégorie $\mathcal{M}$ est connue
à équivalence près quand on connait l'ensemble $\Sigma = \Sigma(\mathcal{M})$,
et l'application
\[
\sigma \longmapsto (Z_{\sigma}, \xi_{\sigma} \in \mathrm{Br}(Z_{\sigma}))
\]
associant à toute classe d'objet simple $M$ de $S$ \struck{la classe de}
\add{le centre $Z_{\sigma}$} $\operatorname{End}(M)$
\note{La machine écrit $S$ pour $\Sigma$ ; on laisse comme sur la page. La
phrase se poursuit page 62.}

\page{62}
et la classe de $\operatorname{End}(M)$ dans $\mathrm{Br}(Z_{\sigma})$ (le
couple $(Z_{\sigma}, \xi_{\sigma})$ étant défini à isomorphisme près). Sans
supposer $\mathcal{M}$ semi-simple, on peut \add{en un certain sens} dire que
la plus grande sous-catégorie semi-simple (engendrée par les objets
semi-simples) est connue à équivalence près par la connaissance précédente,
qui implique donc la connaissance de
$K(\mathcal{M}) = \underline{\mathbb{Z}}^{(\Sigma)}$ \add{\ill{}}, et sa
variance pour les changements de base par $k'$ variable sur $k$. De façon
précise, $\Sigma' = \Sigma(\mathcal{M}_{k'})$ en tant qu'ensemble sur
$\Sigma$ correspond à l'ensemble \struck{des corps} dont la fibre en $\sigma$
est l'ensemble des corps composants de $Z_{\sigma} \otimes_{k} k'$, soient
$Z'^{i}_{\sigma}$, l'application
$K(\mathcal{M}) = \underline{\mathbb{Z}}^{(\Sigma)} \to
K(\mathcal{M}_{k'}) = \underline{\mathbb{Z}}^{(\Sigma')}$ est donné par
\[
\sigma \longmapsto n_{i}(\sigma, i) ,
\]
où $n_{i}$ est donné par la règle suivante : si $d$ est \struck{le degré de}
l'ordre de $\xi_{\sigma} \in \mathrm{Br}(Z_{\sigma})$, celui de son image dans
$\mathrm{Br}(Z'^{i}_{\sigma})$ est un \struck{\ill{}} diviseur $d'_{i}$ de
$d$, et on a $n_{i} = d/d_{i}$. (À généraliser au cas où les $Z_{\sigma}$ ne
sont pas nécessairement séparables). Si $k'$ est une \struck{\ill{}} extension
alg.\ close de $k$, l'ensemble $\Sigma'$ s'identifie à l'ensemble somme des
\emph{fibres} en $k'$ des $\operatorname{Spec}(Z_{\sigma})$ …

b) \quad Supposons maintenant $\mathcal{M}$ \struck{semi-tannakienne} munie
d'un produit tensoriel à contraintes d'associativité et commutativité, avec
$1$. On désigne sous le nom de « caractères numériques » de $\Sigma$ les
données suivantes :

\marginal{préciser cette notion, en faisant d'un « caractère numérique » une
structure susceptible \uncertain{de variation} fonctorielle …}

\begin{enumerate}
\item l'ensemble (à bijection près) \add{$\Sigma$} des \struck{\ill{}} classes
d'isomorphie d'objets simples de $\mathcal{M}$.

\item La donnée de la multiplication de
$K(\mathcal{M}) = \underline{\mathbb{Z}}^{(\Sigma)}$ ; ce qui revient à la
donnée pour tout $\sigma, \tau \in \Sigma$ du produit
\[
\sigma\tau = \sum_{\rho} c^{\rho}_{\sigma,\tau}\, \rho
\qquad (c^{\rho}_{\sigma,\tau} \in \underline{\mathbb{N}}),
\]
\add{i.e. du système d'entiers naturels $\sigma^{\rho}_{\sigma,\tau}$.}
\note{La dernière ligne de la page est de sa main, et elle écrit
$\sigma^{\rho}_{\sigma,\tau}$ là où la formule tapée deux lignes plus haut
porte $c^{\rho}_{\sigma,\tau}$ ; on laisse les deux notations comme elles
sont.}
\end{enumerate}

\page{63}
\begin{enumerate}
\setcounter{enumi}{2}
\item La donnée de l'application
\[
\sigma \longmapsto (Z_{\sigma}, \xi_{\sigma}) , \qquad
\xi_{\sigma} \in \mathrm{Br}(Z_{\sigma})
\]
envisagée dans a).

\item La donnée de la fonction
\[
\text{rang} \;:\; \mathcal{M} \longrightarrow \underline{\mathbb{Z}}
\]
lorsque $\mathcal{M}$ est supposée \struck{\ill{}} tannakienne.
\struck{\ill{}}
\end{enumerate}

On notera que la donnée de 1), 2) équivaut à la donnée de l'anneau commutatif
$K = K(\mathcal{M})$, avec sa « partie effective »
\[
K^{\mathrm{eff}} = \text{sous-groupe de } K(\mathcal{M})
\text{ engendré par les } c\ell(M),\ M \in \mathrm{Ob} ,
\]
\[
= \underline{\mathbb{N}}^{(\Sigma)}
\]
(NB $\Sigma$ se récupère à partir de \struck{ces} \add{la} donnée
\add{$K^{\mathrm{eff}} \subset K$} comme l'ensemble des éléments non nuls
minimaux de $K^{\mathrm{eff}}$). La donnée de 3) permet alors de trouver,
comme explicité dans a), la variance fonctorielle de l'anneau
$K(\mathcal{M}_{k'})$ par rapport à $k'$. Enfin la donnée de 4) permet en
principe de retrouver les polynômes caractéristiques (en particulier, la
trace et le déterminant) d'un endomorphisme d'un objet \struck{\ill{}}
(semi-simple ?) $M$ de $E$, en termes de polynômes caractéristiques réduits
(resp.\ trace réduite, norme réduite) dans les algèbres d'Azumaya
\add{$\operatorname{End}(M_{i})$} correspondant aux composantes isotypiques
$M_{i}$ de $M$. Pour $M$ isotypique de rang $n$, on aura en effet, si
$K = \operatorname{End}(M)$ a comme centre $Z$ de rang $r$ sur $k$, donc
$n = n'r$, et $K$ de rang $d^{2}$ sur $Z$, avec $d \mid n'$ (NB $M$ de rang
$n'$ sur $Z$) :
\begin{align*}
\operatorname{det}_{M} f &= \operatorname{N}_{Z/k}(\operatorname{Nr}_{K/Z}(f))^{n'/d} , \qquad (n'/d = n/dr) \\
\operatorname{Tr}_{M} f &= \operatorname{Tr}_{Z/k}(\operatorname{Trr}_{K/Z}(f))\, \frac{n'}{d} \\
\operatorname{Pol}_{M}(f,t) &= \operatorname{N}_{Z/k}(\operatorname{Polr}_{K/Z}(f,t))^{n'/d}
\end{align*}

\marginal{Signaler que la connaissance des car.\ num.\ de $\mathcal{M}$
implique celle des car.\ numériques de $\mathcal{M}_{k'}$ pour toute
extension $k'$ de $k$, (« fonctoriellement en $k'$ »)}

On fera attention que la connaissance des caractères numériques de la
catégorie tannakienne $\mathcal{M}$ ne permet pas en général de réconstituer
celle-ci à équivalence près, même si $k$ est un corps alg.\ clos de car.\ $0$
(de sorte que la donnée de 3) devient vide) et $\mathcal{M}$ semi-simple

\page{64}
et \add{de plus} même dans le cas particulier où $\mathcal{M}$ est la
catégorie des représentations d'un groupe fini (?\,? à vérifier dans
littérature qu'un groupe fini n'est pas déterminé par la connaissance de la
table de multiplication de \struck{\ill{}} son ensemble de caractères, et les
rangs de ceux-ci …) Comme exception à cette remarque, signalons cependant
le cas où $\mathcal{M}$ est tannakienne à lien diagonalisable (alors $\Sigma$
devient un sous-groupe de $K^{*}$, \add{qui} \struck{dans} détermine le lien,
et la donnée de 3) donne l'homomorphisme
$\Sigma \to \mathrm{Br}(k)$ qui suffit à tout déterminer). Regarder le cas où
le lien est un groupe de t.m.

c) \quad Cas d'un lien qui est de t.m. On suppose connu le caractère
numérique \struck{\ill{}} de $\mathcal{M}$. Prenant le changement de base par
une clôture séparable $k'$ de $k$, on trouve que
$\Sigma' = \Sigma(\mathcal{M}_{k'})$ (canoniquement déduit de la donnée
numérique, avec l'opération de $\pi = \operatorname{Gal}(k'/k)$ dessus) est
muni d'une structure de groupe induite par la multiplication de
$K(\mathcal{M}_{k'}) = \underline{\mathbb{Z}}^{(\Sigma')}$,
\struck{\ill{}} respectée par l'action de $\pi$. Ceci nous donne le lien
$G = D(\Gamma)$ (où $\Gamma$ est le groupe $\Sigma'$ avec l'action de $\pi$).
L'ensemble $\Sigma$ s'identifie à l'ensemble des orbites $\sigma$ de $\pi$
opérant sur $\Gamma$. Si $\sigma$ est une telle orbite, le corps
$Z_{\sigma}$ n'est autre que le corps des invariants
$Z_{\sigma} \subset k'$ du stabilisateur d'un point de l'orbite $\sigma$. Si
$M_{\sigma}$ est un objet de type $\sigma$, on a un homomorphisme
\[
G \longrightarrow \underline{\operatorname{Aut}}(M_{\sigma}) = \underline{K}^{*}_{\sigma}
\]
(les soulignés indiquent qu'on prend des schémas sur $k$ …) qui se
factorise par le centre $Z_{\sigma}$ de $K_{\sigma}$ en
\[
u_{\sigma} : G \longrightarrow \underline{Z}^{*}_{\sigma}
\;\add{\simeq}\; \prod\nolimits_{Z/k} \underline{G}_{m} ,
\]
et on vérifie aussitôt qu'on a
\[
u_{\sigma}(\xi) = \xi_{\sigma} ,
\]
où

\page{65}
\[
\xi \in H^{2}(k_{\mathrm{pl}}, G)
\]
est la classe de la gerbe tannakienne qui définit $\mathcal{M}$. Donc la
connaissance du caractère numérique de $\mathcal{M}$ équivaut à la donnée, en
plus du \struck{\ill{}} groupe de t.m.\ $G$ sur $k$, des classes
$u_{\sigma}(\xi)$ pour tous les homomorphismes $u$ de $G$ dans des tores de
la forme $\prod_{Z/k} \underline{G}_{m}$, où $Z$ est une extension finie de
$k$. \struck{(}On vient de voir en effet que la donnée du caractère numérique
permet de reconstituer ce qui précède. Pour voir l'inverse, notons que
\struck{si} le rang $n(\sigma)$ de $M_{\sigma}$ est déterminé par la formule
\[
n(\sigma) = \operatorname{card}(\sigma)\cdot\operatorname{ord}(\xi_{\sigma})
\]
(où $\operatorname{card}\sigma$ est le cardinal de l'orbite $\sigma$,
$\operatorname{ord}(\xi_{\sigma})$ est l'ordre de $\xi_{\sigma}$ dans
$\mathrm{Br}(Z_{\sigma})$).

La question de savoir si le caractère numérique détermine à équivalence de
catégories tannakiennes près revient donc, pour $G$ fixé, à la question de
savoir si un élément \struck{\ill{}} $\xi \in H^{2}(k_{\mathrm{pl}}, G)$ est
connu quand on connait les $u_{\sigma}(\xi)$ pour $u : G \to \underline{Z}^{*}$
comme dessus. \struck{\ill{}} C'est le cas pour $G$ diagonalisable
\struck{pro-dénombrable} (pro-dénombrable ?) On va examiner deux autres cas.

\emph{Exemple 1}. $G$ est un tore de dimension $1$, donc (s'il n'est pas
trivial) tordu par une extension quadratique $Z$ de $k$. On a donc une suite
exacte
\[
0 \longrightarrow G \xrightarrow{\ u\ } \underline{Z}^{*}
\xrightarrow{\ N_{Z/k}\ } \underline{G}_{m} \longrightarrow 0 ,
\]
d'où on conclut, par la suite exacte de cohomologie et le th.\ 90 que
\[
\operatorname{Ker}\bigl(u : H^{2}(k,G) \longrightarrow H^{2}(k, \underline{Z}^{*})\bigr)
\]
est nul. Donc dans le cas envisagé, le caractère numérique détermine
$\mathcal{M}$ à équivalence près.

\struck{\emph{Exemple 2}. Gardons \ill{} $G$ comme dessus, soit $n$ un
entier $\geq 2$, et soit $G' = \operatorname{Ker} n\cdot\mathrm{id}_{G}$. La
suite exacte}
\[
0 \longrightarrow G' \longrightarrow G \xrightarrow{\ n\ } G \longrightarrow 0
\]
\struck{et le résultat précédent montre que l'ensemble des classes de}
\note{L'Exemple 2 et sa suite exacte, dont la machine n'a tapé ni les flèches
ni les zéros aux bons endroits, sont annulés par de longs traits diagonaux,
et le tapuscrit s'interrompt là : la page 66 ouvre un autre texte.}

\section*{Catégories tannakiennes définies par des cristaux (pages 66 à 80)}

\note{Second tapuscrit, corrigé de la même main. Les paragraphes sont
numérotés de 1 à 9 et courent sans interruption jusqu'à la dernière page du
lot. $W$ désigne un $p$-anneau de corps résiduel $k$ et $K$ son corps des
fractions, $\sigma$ le Frobenius, $\Pi_{p}(k)$ la gerbe des
$\underline{\mathbb{Q}}_{p}$-faisceaux constructibles ; les majuscules
soulignées sont les schémas, comme page 64. La machine n'avait ni le $\Pi$
majuscule ni les accents circonflexes, et la page 71 porte à ce sujet une
note de sa main.}

\page{66}
\emph{Catégories tannakiennes définies par des cristaux.}

1. \quad Soit $k$ un corps de car.\ $p > 0$, qu'on regarde comme algèbre sur
$\underline{\mathbb{Z}}_{p}$, dont l'idéal maximal est muni de puissances
divisées. Cela donne un sens au site cristallin de $k$ (sur
$\underline{\mathbb{Z}}_{p}$, qualifié aussi de « absolu »), et aux Modules
loc.\ libres de type fini \add{(resp.\ de présentation finie)} sur ledit,
qu'on appellera aussi cristaux en modules (localement libres)
\add{resp.\ de prés.\ finie} sur $k$. Ces cristaux forment une
$\otimes$-catégorie $\underline{\mathbb{Z}}_{p}$-linéaire,
\struck{\ill{}} On peut expliciter cette catégorie à l'aide d'un
$p$-anneau $W$ de corps résiduel $k$, comme la catégorie des modules libres
de type fini (resp.\ de type fini) $M$ sur $W$, avec pour tout entier $n$ une
connexion absolue sur $M_{n} = M/p^{n}M$ (relativement à $W_{n}$ sur
$\underline{\mathbb{Z}}/p^{n}\underline{\mathbb{Z}}$), à courbure nulle, ces
connexions se recollant pour $n$ variable. Lorsque $k$ est parfait, $W$ est
l'anneau des vecteurs de Witt construit sur $k$, et dans la description
précédente, on peut oublier les connexions : les catégories envisagées sont
simplement la catégorie des modules libres de t.f (resp.\ des modules de
t.f.) sur $W$. Si $K$ est le corps des fractions de $W$, la catégorie des
\emph{isocristaux} \struck{sur} (=cristaux à isogénie près) sur $k$ est
équivalente à $\operatorname{Modf}(K)$. \struck{\ill{}}

2. \quad La catégorie des cristaux sur $k$ dépend fonctoriellement de $k$,
d'où pour tout cristal la possibilité de définir son transformé par le
changement de base $\sigma_{k} : x \mapsto x^{p} : k \to k$, soit
$M^{\sigma}$. La donnée d'un $M$ avec un homomorphisme
\[
F_{M} : M^{\sigma} \longrightarrow M
\]
s'appelle un \emph{$F$-cristal sur} $k$. On notera que les
$F$-cristaux forment une \add{$\underline{\mathbb{Z}}_{p}$-linéaire}
catégorie \struck{$\otimes$-} dont les $\operatorname{Hom}$ sont des modules
de type fini sur $\underline{\mathbb{Z}}_{p}$ (attention, pas
\add{des modules} sur $W$, même si $k$ est parfait !). Ce n'est pas une
\struck{\ill{}} catégorie tannakienne, même une fois tensorisée par
$\underline{\mathbb{Q}}_{p}$, car il manque

\page{67}
un \emph{$\operatorname{Hom}$} interne. On dit qu'un $F$-cristal $M$ est
\emph{non dégénéré} s'il existe un entier $i \geq 0$ et un
homomorphisme $V : M \to M^{\sigma}$ tel que $VF = p^{i}$. Soit
$\mathrm{FCris}^{\ddagger}(k)$ la catégorie \struck{\ill{}} déduite de la
catégorie des \struck{\ill{}} $F$-cristaux non dégénérés en localisant par
rapport à $p$ ; c'est donc une sous-$\otimes$-catégorie de celle des
$F$-cristaux à isogénie près. On l'appellera la catégorie des
\emph{$F$-isocristaux effectifs}, et on laissera tomber le qualificatif
« iso » s'il n'y a pas de risque de confusion.
\emph{Le $F$-isocristal de Tate inverse} \add{$K(-1)$} est déduit du
$F$-cristal défini par le cristal \emph{unité}, avec
$F_{K(-1)} = p\cdot\mathrm{id}$ \add{$W(-1)$} (cet $F$-cristal est évidemment
non dégénéré !). On voit alors que $K(-1)$ est faiblement inversible, et on
constate que lorsqu'on ajoute formellement à $\mathrm{FCris}^{+}(k)$
\struck{\ill{}} un inverse $K(1)$ de $K(-1)$ alors on trouve une catégorie
tannakienne sur $\underline{\mathbb{Q}}_{p}$, ayant un foncteur fibre évident
sur $K$ ; cette catégorie est notée $\mathrm{FCriso}(k)$, et appelée
\emph{catégorie des $F$-isocristaux} (ou simplement, s'il n'y a
\struck{pas} pas de risque de confusion, des $F$-cristaux) sur $k$. $K(1)$
est le $F$-cristal de Tate, et on désigne par $M \mapsto M(n)$ le foncteur
de tensorisation par sa puissance tensorielle $n$.ème. Si $n = -m$,
$m \geq 0$, ce foncteur est induit par le foncteur sur les $F$-cristaux
effectifs (sans \struck{isogénie}) consistant à multiplier $F$ par $p^{m}$.

3. \quad Lorsque $k$ est parfait, désignant par $\sigma$ l'automorphisme de
Frobenius de $W = W(k)$, la donnée d'un $F$-cristal revient à la donnée d'un
module de type fini $M$ sur $W$ muni d'un homomorphisme
\[
F_{M} : M^{\sigma} = M \otimes_{W} (W, \sigma) \longrightarrow M ,
\]
et le $F$-cristal est non dégénéré sss l'homomorphisme précédent est une
isogénie i.e.\ est de rang égal à rang $M$, ou encore sss il induit un
isomorphisme après tensorisation par $K$. La catégorie
$\mathrm{FCris}^{+}(k)$ s'identifie alors à la catégorie des vectoriels de
dimension finie

\page{68}
$E$ sur $K$, munis d'un \emph{automorphisme} $\sigma$-linéaire $F_{E}$,
qu'on a intérêt à interpréter comme un isomorphisme $K$-linéaire
\[
F_{E} : E^{\sigma} = E \otimes_{K} (K, \sigma) \longrightarrow E .
\]
Les opérations tensorielles sont définies par transport de structure (cf.\ cas
d'un groupe, ici $\underline{\mathbb{Z}}$, opérant sur un corps …). Les
$F$-isocristaux effectifs s'interprètent alors comme les $(E, F_{E})$ tels que
$F_{E}$ laisse invariant un réseau, ce qui signifie aussi que l'ensemble des
itérés $F_{E}^{n}$, $n \geq 0$, est \emph{borné} dans le vectoriel
$\operatorname{End}(E)$. Tout $(E, F_{E})$ devient \add{évidemment}
\struck{effectif} après multiplication de $F_{E}$ par une puissance
convenable de $p$, ce qui est une autre façon de dire que la catégorie des
$F$-isocristaux se déduit de celle des $F$-isocristaux effectifs par en
rendant inversible $K(-1)$.

4. \quad La catégorie tannakienne $\mathrm{FCriso}(k)$ est un invariant
arithmétique intéressant attaché à $k$ (fonctoriellement) ; sa
\struck{\ill{}} connaissance équivaut à celle de la gerbe associée (sur
$\underline{\mathbb{Q}}_{p}$), soit $\underline{G}(k)$. Parmi les structures
supplémentaires qu'il convient d'étudier en même temps que cette gerbe, il
faut inclure la structure d'effectivité de la catégorie
$\mathrm{FCriso}(k)$ dans ses représentations linéaires, celle de l'objet de
Tate inverse (définissant un homomorphisme canonique de $\underline{G}(k)$
vers la gerbe triviale de groupe $\underline{G}_{m}$), enfin pour $k$ parfait
tout au moins, la donnée du foncteur fibre naturel sur $K(k)$ (qui est donc
une section de la gerbe $\underline{G}(k)$ sur $K(k)$) ; pour $k$ non
parfait, il conviendrait d'étudier de même la section de
$\underline{G}(k)$ sur le corps des fractions $K$ de n'importe quel
$p$-anneau $W$ de corps résiduel $k$.

L'intérêt de la notion de $F$-cristal \struck{non dégénéré} provient du fait
que l'on a un foncteur « cohomologie cristalline »
\[
X \longmapsto H^{*}_{\mathrm{cris}}(X)
\]

\page{69}
allant de la catégorie des schémas propres et lisses sur $k$ vers les
$F$-cristaux gardués, \struck{\ill{}} \add{\ill{}} qui est additif (pour la
somme des $X$) et compatible avec $\otimes$ (le $\otimes$ des schémas étant le
produit cartésien) de façon compatible avec les associativités et
commutativités (tant pour la somme que pour le produit), en faisant attention
d'utiliser la règle de Koko pour la contrainte de commutativité. De plus,
modulo des vérifications qui n'ont pas encore été faites par Berthelot, mais
qui ne peuvent manquer de sortir prochainement, le foncteur précédent est
même défini en prenant comme morphismes des correspondances algébriques à
équivalence algébrique près, et s'en va vers les $F$-cristaux
\emph{non dégénérés} ; de façon précise, il doit être vrai que pour tout
$i$, il existe un $V_{i} : \struck{\ill{}} H^{i} \to (H^{i})^{\sigma}$ tel que
$F V_{i} = p^{i}\,\mathrm{id}$, $V_{i} F = p^{i}\,\mathrm{id}$. (C'est en tous
cas vrai si $X$ se remonte en car.\ $0$) Modulo cette vérification, on trouve
donc un foncteur
\[
H^{*}_{\mathrm{cris}} : \text{schémas propres lisses sur } k \longrightarrow
\add{\mathrm{Grad}}\,(\mathrm{FCriso}^{+}(k)) ,
\]
compatible à tout (un « foncteur cohomologique » dans la terminologie de
l'exposé de Kleiman $[\ ]$). Bien entendu, on a
$H^{2}_{\mathrm{cris}}(\mathbb{P}^{1}) = W(-1)$, plus généralement on devra
avoir (quand Berthelot aura bien travaillé)
\[
H^{2d}_{\mathrm{cris}}(X) \simeq W(-d)
\]
si $X$ est géométriquement connexe et de dimension $d$, (isomorphisme
canonique, compatible à la multiplication …). En admettant le yoga des
motifs (on a besoin ici de la structure graduée de la catégorie des motifs
sur $k$, pour pouvoir canuler par la règle de Koko la contrainte de
commutativité naturelle sur les motifs, provenant de l'isomorphisme de
commutativité $X \times Y \simeq Y \times X$ sur les schémas …) on aura
un homomorphisme de gerbes canonique
\[
\prod(k)_{\underline{\mathbb{Q}}_{p}} \longleftarrow \underline{G}(k) ,
\qquad \add{\text{associé au } \otimes\text{-foncteur }
(\mathrm{Mot}(k) \longrightarrow \mathrm{FCriso}(k))}
\]
où $\prod(k)$ est la gerbe associée à la catégorie tannakienne (sur
$\underline{\mathbb{Q}}$)

\page{70}
des motifs \add{(iso) (semi-simples)} sur $k$, et l'indice
$\underline{\mathbb{Q}}_{p}$ désigne \struck{l'effet du} changement de base
par $\underline{\mathbb{Q}} \to \underline{\mathbb{Q}}_{p}$. Cet
homomorphisme est compatible à l'augmentation « de Tate » et aux structures
d'effectivité. Modulo du travail à faire de Berthelot, pour
$k = \underline{\mathbb{F}}_{q}$ un corps fini à $q$ éléments, la
connaissance de l'image \struck{\ill{}} $H^{*}(X)$ dans
$\mathrm{FCriso}(k)$ permet de calculer la fonction $\zeta_{X}$ de $X$, par
la formule qu'on devine et que je me dispense d'écrire. (D'ailleurs, modulo
les conjectures de Weil, la connaissance de la fonction $\zeta_{X}$ permet,
grâce à cette formule, de réconstituer $H^{*}_{\mathrm{cris}}(X)$ à
isomorphisme près dans $\mathrm{FCriso}(k)$, du moins virtuellement (i.e.\
l'image de chaque $H^{i}_{\mathrm{cris}}(X)$ dans l'anneau
$K(\mathrm{FCriso}(k))$) si on ne désire pas admettre que \struck{\ill{}} les
$H^{i}$ sont nécessairement des éléments semi-simples de
$\mathrm{FCriso}(k)$ …)

5. \quad $F$-\emph{cristaux de \struck{niveau} pente nulle}. On définira
plus loin la \emph{pente} d'un $F$-cristal « homogène ». Ici, nous
allons introduire directement les $F$-cristaux de pente nulle. Pour ceci,
pour tout faisceau $p$-adique constructible \add{$V$} sur $k$ (correspondant
donc à une \struck{représentation} \add{opération} continue de
$\pi = \operatorname{Gal}(\bar{k}/k)$ dans un module de type fini sur
$\underline{\mathbb{Z}}_{p}$) on lui associe de façon évidente un $F$-cristal
$M$, noté parfois (par abus de notations)
\[
M = V \otimes_{\underline{\mathbb{Z}}_{p}} W ,
\]
(même s'il n'y a pas de $W$ choisi, dans le cas $k$ imparfait !). On aura :

\emph{Proposition 5.1}. Le foncteur
$M \mapsto V \otimes_{\underline{\mathbb{Z}}_{p}} W$ est une équivalence de
la catégorie des faisceaux $p$-adiques constructibles sur $k$ vers la
sous-catégorie pleine de la catégorie des $F$-cristaux \struck{sur} sur $k$,
formée des $F$-cristaux $M$ tels que
$F_{M} : M^{\sigma} \to M$ soit un \emph{isomorphisme}.

\emph{Corollaire 5.1}. Le foncteur précédent induit un foncteur
pleinement fidèle (compatible avec les $\otimes$-structures)
\add{$i$} de la catégorie des

\page{71}
$\underline{\mathbb{Q}}_{p}$-faisceaux constructibles sur $k$ vers la
catégorie $\mathrm{FCriso}(k)$. Pour un objet \struck{\ill{}} $(E, F_{E})$ de
cette dernière, les conditions suivantes sont équivalentes
\add{(si $k$ parfait)} : a) $(E, F_{E})$ appartient à l'image essentielle du
foncteur précédent ; \quad b) il existe un réseau $M \subset E$ stable par
$F_{E}$ \struck{et} \add{et} par $F_{E}^{-1}$ (i.e.\ tel que $F_{E}$ induise
un \emph{automorphisme} de $M$) ; \quad c) la famille des itérés
$F^{n}$, $n \in \underline{\mathbb{Z}}$, est bornée ; \quad d) $E$ et
$\check{E}$ sont effectifs.

Les $F$-isocristaux en question sont appelés « trivaux » ou
« \emph{de pente nulle} ». On trouve ainsi un
\emph{épimorphisme de gerbes}
\[
\underline{G}(k) \longrightarrow \underline{\Pi}_{p}(k) ,
\]

\marginal{$\Pi$ majuscule ?}

où $\underline{\Pi}_{p}(k)$ est la gerbe sur $\underline{\mathbb{Q}}_{p}$
\struck{\ill{}} associée à la catégorie tannakienne des
$\underline{\mathbb{Q}}_{p}$-faisceaux constructibles, équivalente
\struck{\ill{}} (une fois choisie une clôture algébrique $\bar{k}$ de $k$) à
la gerbe triviale définie par l'enveloppe algébrique \struck{du} $p$-adique
du groupe pro-fini $\pi = \operatorname{Gal}(\bar{k}/k)$. Tenant compte de
l'homomorphisme d'augmentation de Tate, on trouve même un homomorphisme
\[
\underline{G}(k) \longrightarrow \underline{\Pi}_{p}(k) \times
[\underline{G}_{m}]_{\underline{\mathbb{Q}}_{p}}
\]
où $[\underline{G}_{m}]$ est la gerbe triviale définie par
$\underline{G}_{m}$. On vérifie facilement que c'est encore un épimorphisme
(i.e.\ que le foncteur
\[
(V_{i})_{c} \longmapsto \sum V_{i} \otimes \underline{\mathbb{Q}}_{p}(i)
\]
allant de la catégorie des \struck{\ill{}}
$\underline{\mathbb{Q}}_{p}$-faisceaux constructibles \add{gradués} sur $k$
vers $\mathrm{FCriso}(k)$ est pleinement fidèle). Les $F$-isocristaux de la
forme précédente pourraient être appelés les $F$-isocristaux
\emph{de niveau zéro}.

6. \quad Considérons maintenant un homomorphisme de corps \struck{\ill{}}
\[
k \longrightarrow k' ,
\]
d'où un homomorphisme de catégories tannakiennes sur
$\underline{\mathbb{Q}}_{p}$
\[
\mathrm{FCriso}(k) \longrightarrow \mathrm{FCriso}(k') ,
\]
et un homomorphisme de gerbes tannakiennes sur
$\underline{\mathbb{Q}}_{p}$ en sens inverse

\page{72}
\[
\underline{G}(k') \longrightarrow \underline{G}(k) .
\]
\add{(mod isomorphisme canonique)} Ce dernier rend commutatif le diagramme

\begin{tikzcd}[row sep=small, nodes={font=\scriptsize}]
\underline{G}(k) \arrow[r] \arrow[d] & \underline{G}(k') \arrow[d] \\
\underline{\Pi}(k') \times [\underline{G}_{m}]_{\underline{\mathbb{Q}}_{p}} \arrow[r] & \underline{\Pi}_{p}(k) \times [\underline{G}_{m}]_{\underline{\mathbb{Q}}_{p}}
\end{tikzcd}

(i.e.\ on a, sur les catégories tannakiennes \struck{\ill{}} correspondantes,
un diagramme commutatif à isomorphisme \add{can.} près).

Oubliant dans ce dernier les facteurs $[\underline{G}_{m}]$, on voit donc que
$\underline{G}(k') \to \underline{G}(k)$ se factorise par la
sous-\struck{catégorie} gerbe de $\underline{G}(k)$ correspondant au
sous-groupe fermé de $\pi_{1}(k)$ image de $\pi_{1}(k')$ …\ A prouver :

\emph{Proposition 6.1}. Si $k'$ est une extension \struck{finie} étale de
$k$, alors l'homomorphisme de gerbes
$\underline{G}(k') \to \underline{G}(k)$ est un
\emph{monomorphisme} dont l'image est exactement la sous-gerbe
« ouverte » définie ci-dessus. En particulier, si $k'$ est une extension
finie galoisienne de $k$, \struck{\ill{}} de groupe de Galois $g$, et on une
suite exacte canonique de gerbes
\[
0 \longrightarrow \underline{G}(k') \longrightarrow \underline{G}(k)
\longrightarrow [g]_{\underline{\mathbb{Q}}_{p}} \longrightarrow 0 .
\]

Cet énoncé devrait être formellement équivalent au suivant :
$\operatorname{Spec}(k') \to \operatorname{Spec}(k)$ est un morphisme de
descente effective pour la catégorie fibrée des $\mathrm{FCriso}(L)$ sur des
composés finis \add{$L$} de corps de car.\ $p$. La démonstration de ce
dernier fait est $\pm$ triviale.

\emph{Question 6.2}. A-t-on encore un énoncé analogue lorsqu'on suppose
que $k'$ est une extension galoisienne de $k$ qui peut-être infinie ? Le cas
le plus intéressant serait celui où $k'$ est la clôture séparable de $k$, de
sorte qu'on aurait un dévissage de $\underline{G}(k)$ en
\[
\add{(6.2.1)} \qquad 0 \longrightarrow \underline{G}(\bar{k}) \longrightarrow
\underline{G}(k) \longrightarrow \underline{\Pi}_{p}(k) \longrightarrow 0 ;
\]
lorsque $k$ est parfait i.e.\ $\bar{k}$ alg.\ clos, nous verrons que
$\underline{G}(\bar{k})$ ne dépend pas essentiellement de $k$. À

\page{73}
7. \quad Nous verrons ci-dessus que la réponse à la question précédente est
affirmative lorsque $k$ est un corps fini, en particulier si c'est le corps
premier $\underline{\mathbb{F}}_{p}$. On trouve alors, grâce à cette
structure d'extension sur $\underline{G}(\underline{\mathbb{F}}_{p})$, un
homomorphisme canonique de Gerbes (pour tout $k$)
\[
\underline{G}(k) \simeq \underline{G}(\underline{\mathbb{F}}_{p})
\times_{\underline{\Pi}_{p}(\underline{\mathbb{F}}_{p})}
\underline{\Pi}_{p}(k) ,
\]
et pour $k$ parfait, une réponse affirmative à la question précédente revient
à l'assertion que l'homomorphisme précédent est une équivalence de Gerbes.
Comme on connait parfaitement bien la Gerbe
$\underline{G}(\underline{\mathbb{F}}_{p})$ et son homomorphisme dans
$\underline{\Pi}_{p}(\underline{\mathbb{F}}_{p})$ (cf plus bas) cela
donnerait donc la structure de $\underline{G}(k)$ en termes de la
connaissance de \struck{\ill{}} $\underline{\Pi}_{p}(k)$ (connu quand on
connait $\pi_{1}(k) = \operatorname{Gal}(\bar{k}/k)$) et de l'homomorphisme
$\underline{\Pi}_{p}(k) \to \underline{\Pi}_{p}(\underline{\mathbb{F}}_{p})$
(connu quand on connait
$\pi_{1}(k) \to \pi_{1}(\underline{\mathbb{F}}_{p}) \to \hat{Z}$). Cela
résoudrait donc la question de la structure de $\underline{G}(k)$ pour $k$
parfait. On notera que l'objet de Tate de $\mathrm{FCriso}(k)$ provient de
celui de $\mathrm{FCriso}(\underline{\mathbb{F}}_{p})$ et sera donc également
déterminé. Le foncteur fibre canonique sur $K$ \add{$= K(k)$} correspond
alors à un foncteur fibre de
$\mathrm{FCriso}(\underline{\mathbb{F}}_{p})$ sur $K$, un foncteur fibre de
\struck{\ill{}} \add{la catégorie des $\underline{\mathbb{Q}}_{p}$-faisceaux
constr.\ sur $k$} sur $K$, et un isomorphisme entre les « restrictions » de
ces foncteurs à la catégorie des $\underline{\mathbb{Q}}_{p}$-faisceaux
constructibles sur $\underline{\mathbb{F}}_{p}$ ; or le premier foncteur
n'est autre que \add{celui déduit par} \struck{l'}extension des scalaires
$\underline{\mathbb{Q}}_{p} = K(\underline{\mathbb{F}}_{p}) \to K = K(k)$ du
foncteur fibre analogue sur $\mathrm{FCriso}(\underline{\mathbb{F}}_{p})$, le
deuxième est le foncteur
$V \mapsto V \otimes_{\underline{\mathbb{Q}}_{p}} K$, et l'isomorphisme
\struck{\ill{}} entre les restrictions est l'isomorphisme canonique évident
\add{(compatibilité de l'anneau de \ill{} des corps de base)}. Donc la
structure supplémentaire de $\underline{G}(k)$ (savoir sa section
\struck{sur} canonique sur $K$) est déterminée par la connaissance de la
section de $\underline{\Pi}_{p}(k)$ sur $K$ (i.e.\ le foncteur précédent
$V \mapsto V \otimes_{\underline{\mathbb{Q}}} K$) et l'isomorphisme de
l'image de ladite section dans
$\underline{\Pi}_{p}(\underline{\mathbb{F}}_{p})$ avec la section provenant
de \struck{\ill{}} la section canonique de
$\underline{\Pi}_{p}(\underline{\mathbb{F}}_{p})$ sur
$\underline{\mathbb{Q}}_{p}$ ; ce sont là des données qui ne sont évidemment
pas contenues dans celle du groupe

\marginal{\textbf{N.B.} On a donc, en comparant le foncteur
« fibre \struck{\ill{}} en $\bar{k}$ » sur $\underline{\mathbb{Q}}_{p}$, et
le foncteur précédent, un torseur canonique sur $K$ de groupe
\uncertain{de lien} $\underline{\Pi}_{p}(k)_{K}$ \struck{\ill{}}}

\page{74}
pro-fini « abstrait » $\pi_{1}(k)$, mais on se convainc facilement qu'elles
sont contenues dans la donnée de ce groupe, $+$ son opération sur la clôture
algébrique $\bar{K}$ de $K$ (ou simplement, sur sa clôture non ramifiée sur
$W$ …). $-$ Il resterait enfin à déterminer la structure d'effectivité
dans la catégorie $\mathrm{FCriso}(k)$ \struck{définie par la récette
précédente} \add{\ill{}} \struck{\ill{}} Or il est immédiat que l'effectivité
d'un $F$-isocristal est invariante par changement de base, en particulier par
le changement de base $k \to \bar{k}$, ce qui ramène au cas d'un $k$
algébriquement clos, ou encore au cas de
$\underline{\mathbb{F}}_{p}$. \struck{Dans les termes de} la terminologie
introduite plus bas, on trouvera donc que les $F$-isocristaux effectifs sont
ceux qui sont \emph{à pente} $\geq 0$ (i.e.\ dont toutes les composantes
isopentiques \struck{\ill{}} non nulles sont de pente $> 0$).

\marginal{On a un hom.\ (une fois choisi $\bar{k}$, d'où une trivialisation
de lien $\underline{\Pi}_{p}(k)$) $\pi_{1}(k)_{K} \to G'(K)$, d'où un
foncteur (torseurs sous $\pi_{1}(k)_{K}$) $\to$ torseurs sous $G'$, et on
prend l'image d'un torseur \struck{sous} $\pi_{1}(k)_{K}$, \uncertain{défini}
par \ill{}, interprété en $\pi_{1}(k)$ comme \ill{} de l'extension \ill{} non
ramifiée \ill{} de $k$.}

8. \quad Il reste donc à étudier la gerbe
$\underline{G}(\underline{\mathbb{F}}_{p})$, avec sa \struck{\ill{}}
\struck{\ill{}} section canonique \emph{sur $\underline{\mathbb{Q}}_{p}$},
sa structure d'extension, et sa structure d'effectivité. \add{d.1.} Tout
d'abord, la donnée du foncteur fibre canonique permet d'identifier
$\underline{G}(\underline{\mathbb{F}}_{p})$ à la gerbe triviale définie par
le lien $\underline{G}$ de cette gerbe. Comme
$\mathrm{FCriso}(\underline{\mathbb{F}}_{p})$ avec son foncteur fibre
s'identifie à la catégorie des automorphismes \add{$F_{E}$} de vectoriels de
dimension finie \add{$E$} sur $\underline{\mathbb{Q}}_{p}$, avec le foncteur
« oubli de $F_{E}$ », on voit que le lien de la gerbe est l'enveloppe
$\underline{\mathbb{Q}}_{p}$-\add{pro}algébrique du groupe discret
$\underline{\mathbb{Z}}$. Soit $G$ ce groupe affine sur
$\underline{\mathbb{Q}}_{p}$. Il est évidemment commutatif, donc produit d'un
groupe vectoriel $G_{u}$ par un groupe de t.m.\ $G_{r}$. Ce dernier est
l'enveloppe \add{de type} multiplicative du groupe $\underline{\mathbb{Z}}$,
et il s'ensuit aussitôt que c'est le groupe de t.m.\ dont le groupe des
caractères sur $\bar{\underline{\mathbb{Q}}}_{p}$ est donné par
\[
\Gamma = \bar{\underline{\mathbb{Q}}}^{*}_{p} ,
\]
(isomorphisme compatible avec l'action de
$\operatorname{Gal}(\bar{\underline{\mathbb{Q}}}_{p}/\underline{\mathbb{Q}}_{p})$).
D'autre

\page{75}
\marginal{une fois \struck{choisi} \uncertain{une clôture} de
$\bar{\underline{\mathbb{F}}}_{p}$}
part, la gerbe
$\underline{\Pi}_{p}(\underline{\mathbb{F}}_{p})$, grâce au \struck{\ill{}}
foncteur fibre évident « fibre en $\bar{\underline{\mathbb{F}}}_{p}$ » sur la
catégorie des $\underline{\mathbb{Q}}_{p}$-faisceaux constructibles sur
$\underline{\mathbb{F}}_{p}$, s'identifie à la \struck{\ill{}}
\add{gerbe triviale définie} par l'enveloppe
$\underline{\mathbb{Q}}_{p}$-algébrique du groupe profini
$\pi_{1}(\underline{\mathbb{F}}_{p}) = \hat{Z}$, \struck{\ill{}}

Un faisceau constructible sur $\underline{\mathbb{F}}_{p}$ s'identifie à un
vectoriel de dimension finie $E$ sur $\underline{\mathbb{Q}}_{p}$, muni d'un
automorphisme $F_{E}$ tel que l'\struck{\ill{}} opération
$n \mapsto F_{E}^{n}$ de $\underline{\mathbb{Z}}$ sur $E$ se prolonge par
continuité en une opération de $\hat{Z}$. Identifiant \add{(par un foncteur
$q$, distinct de celui de d.2.)} la catégorie de ces $(E, F_{E})$ à une
sous-catégorie pleine de la catégorie \struck{\ill{}}
$\mathrm{FCriso}(\underline{\mathbb{Q}}_{p})$ des couples $(E, F_{E})$, où
$F_{E}$ est un automorphisme \emph{quelconque}, on voit que le lien
\struck{\ill{}} $\underline{\Pi}_{p}(\underline{\mathbb{F}}_{p})$ de cette
catégorie s'identifie à un quotient du lien $G$ de
$\mathrm{FCriso}(\underline{\mathbb{F}}_{p})$ déterminé plus haut, et on voit
aussitôt que ce quotient est de la forme $G_{u} \times G'_{r}$, où $G'_{r}$
est \struck{le} \add{$G'$ de} quotient de $G_{r}$ (ceci revient à dire que si
on décompose $F_{E}$ en produit de sa composante unipotente par sa composante
semi-simple, alors la condition que $F_{E}$ définisse une opération de
$\hat{Z}$ ne dépend que de la composante semi-simple) ; de façon précise,
$G'_{r}$ est le groupe quotient de $G_{r}$ défini par le sous-groupe du
groupe des caractères $\bar{\underline{\mathbb{Q}}}^{*}_{p}$ de $G_{r}$ formé
des \emph{unités} $p$-adiques : cela signifie en effet que $F_{E}$
définit une opération de $\hat{Z}$ (i.e.\ définit un $F$-isocristal de pente
nulle) sss les valeurs propres de $F_{E}$ sont des unités, ce qui résulte en
effet aussitôt du critère donné plus haut des $F$-isocristaux de pente nulle
(ou se vérifie directement de façon immédiate ; la référence au résultat sur
les fourbis de pente nulle devrait se justifier par un grain de sel, car
l'inclusion utilisée ici des $\underline{\mathbb{Q}}_{p}$-faisceaux
constructibles sur $\underline{\mathbb{F}}_{p}$ dans les $F$-isocristaux
n'est pas celle utilisée \struck{plus haut} \add{dans d.2}, cf.\ plus bas). On
trouve ainsi une structure d'extension sur
$G \add{=} \operatorname{Lien}(\underline{G}(\underline{\mathbb{F}}_{p}))$ :
\[
0 \longrightarrow H \longrightarrow G \longrightarrow G' \longrightarrow 0 ,
\]
où $H$ est le groupe de t.m.\ dont le groupe des caractères est le

\page{76}
quotient de $\bar{\underline{\mathbb{Q}}}^{*}$ par le sous-groupe des unités,
donc canoniquement isomorphe à $\underline{\mathbb{Q}}$, par l'isomorphisme
\[
\bar{\underline{\mathbb{Q}}}^{*}/\text{unités} \longrightarrow
\underline{\mathbb{Q}}
\]
donnés par la valuation $p$-adique $v_{p}$, normalisée par la condition
\[
v_{p}(p) = 1 .
\]
L'opération de \struck{\ill{}}
$\operatorname{Gal}(\bar{\underline{\mathbb{Q}}}_{p}/\underline{\mathbb{Q}}_{p})$
sur ce groupe quotient $Q$ est triviale, de sorte que $H$ est
\emph{diagonalisable} de groupe des caractères $\underline{\mathbb{Q}}$ :
\[
H = D(\underline{\mathbb{Q}}) .
\]
L'augmentation de Tate de $G$ correspond à l'homomorphisme
\[
\underline{\mathbb{Z}} \longrightarrow \bar{\underline{\mathbb{Q}}}^{*} ,
\qquad 1 \longmapsto p ,
\]
et l'augmentation induite sur $H$ correspond à l'\emph{inclusion}
canonique
\[
\underline{\mathbb{Z}} \hookrightarrow \underline{\mathbb{Q}}
\]
(grâce à notre choix de la valuation \emph{normalisée} comme indiqué).

8.2. \quad Il faut maintenant expliciter l'homomorphisme de gerbes
\[
\add{(8.2.1)} \qquad i : \underline{G}(\underline{\mathbb{F}}_{p})
\longrightarrow \underline{\Pi}_{p}(\underline{\mathbb{F}}_{p})
\]
correspondant à l'inclusion
\[
i : \underline{\mathbb{Q}}_{p}\text{-faisceaux constructibles}
\hookrightarrow \mathrm{FCriso}(\underline{\mathbb{F}}_{p})
\]
de 5.2. On a vu que l'image essentielle est exactement égale à la
sous-catégorie \add{pleine} ($\underline{\mathbb{Q}}_{p}$-faisceaux
constructibles), \struck{\ill{}} de sorte que $i$ s'identifie à une
\struck{\ill{}} $\otimes$-équivalence de la catégorie des
$\underline{\mathbb{Q}}_{p}$-faisceaux constructibles avec elle-même. On
constate alors que l'homomorphisme correspondant sur les liens
\struck{\ill{}} $G' \to G'$ est l'identité, \struck{\ill{}} de sorte que $i$
est défini par un torseur sous $G'$, et que ce torseur sous $G'$ peut se
décrire ainsi (aurait pu passer à la fin de 7, dans le cas d'un $k$
quelconque …). On a un homomorphisme canonique
\[
\add{(8.2.2)} \qquad \pi_{1}(\underline{\mathbb{F}}_{p}) = \hat{Z}
\longrightarrow G'(\underline{\mathbb{Q}}_{p}) , \qquad \text{i.e.}\quad
(\pi_{1}(\underline{\mathbb{F}}_{p}))_{\underline{\mathbb{Q}}_{p}}
\longrightarrow G'
\]
d'où un \struck{homomorphi} foncteur
\[
\text{torseurs sous } (\hat{Z})_{\underline{\mathbb{Q}}_{p}}
\longrightarrow \text{torseurs sous } G'
\]

\page{77}
d'autre part, on a un torseur évident $P_{o}$ sur
$\underline{\mathbb{Q}}_{p}$ de groupe $(\hat{Z})$, correspondant à
l'extension non ramifiée maximale de $\underline{\mathbb{Q}}_{p}$ ; \struck{on}
prend l'image \add{$P$} de ce torseur : c'est lui.

8.3. \quad Soit $\underline{H}$ la gerbe noyau de (8.2.1), qui a donc comme
lien le groupe diagonalisable $H$ de 8.1. D'après le sorite général la classe
\[
\xi \in H^{2}(\underline{\mathbb{Q}}_{p}, \struck{\ill{}}\ H)
= \operatorname{Hom}(\underline{\mathbb{Q}}, \mathrm{Br}(\underline{\mathbb{Q}}_{p}))
= \operatorname{Hom}(\underline{\mathbb{Q}}, \underline{\mathbb{Q}}/\underline{\mathbb{Z}})
\]
de cette gerbe est donnée par
\[
\xi = \partial(\eta)
\]
où
\[
\eta \in H^{1}(\underline{\mathbb{Q}}_{p}, G')
\]
est la classe du torseur $P$. Pour faire le calcul, notons qu'on a un
\add{quotient} \struck{sous-groupe} $H'$ de $G_{r}$, correspondant au
sous-groupe de $\bar{\underline{\mathbb{Q}}}^{*}_{p}$ \struck{formé} engendré
par les racines de l'unité (= éléments de torsion) et toutes les racines de
$p$, $p^{1/n}$ (ce groupe est bien invariant par
$\operatorname{Gal}(\bar{\underline{\mathbb{Q}}}_{p}/\underline{\mathbb{Q}}_{p})$).
On a alors un homomorphisme d'extensions
\[
\begin{array}{ccccccccc}
0 &\longrightarrow& H &\longrightarrow& G &\longrightarrow& G' &\longrightarrow& 0 \\
 && \big\| && \downarrow && \downarrow && \\
0 &\longrightarrow& H &\longrightarrow& H' &\longrightarrow& (\hat{\underline{\mathbb{Z}}})_{\underline{\mathbb{Q}}_{p}} &\longrightarrow& 0 ,
\end{array}
\]
\struck{qui montre qu'on} \struck{\ill{}} a
\[
\xi = \partial(\eta_{o}) ,
\]
où
\[
\eta_{o} \in H^{1}(\underline{\mathbb{Q}}_{p}, \hat{\underline{\mathbb{Z}}})
\]
est l'image de $\eta$ par $G' \to \hat{\underline{\mathbb{Z}}}$. On constate
que le composé de (8.2.2) avec l'homomorphisme précédent est l'identité de
$\hat{\underline{\mathbb{Z}}}$, donc $\eta_{o}$ est la classe du torseur
canonique $P_{o}$ de la fin de 8.2. À l'aide de ceci, il faut vérifier que la
classe $\xi$ correspond (par l'isomorphisme
\[
\mathrm{Br}(\underline{\mathbb{Q}}_{p}) \simeq
\underline{\mathbb{Q}}/\underline{\mathbb{Z}}
\]
du corps de classe local, à l'homomorphisme canonique
$\underline{\mathbb{Q}} \to \underline{\mathbb{Q}}/\underline{\mathbb{Z}}$

\marginal{\textbf{N.B.} Ce diagramme est induit par un diagramme de gerbes
(valable pour un corps $k$ parfait quelconque) où $\underline{H}'$ = gerbe
correspondant à la sous-catégorie de
$\mathrm{FCriso}(\underline{\mathbb{F}}_{p})$ formée des objets qui
deviennent isomorphes par une extension \emph{finie} $k'$ de
$k = \underline{\mathbb{F}}_{p}$ : une somme d'objets $E_{r,s}$ ; et
$\underline{H}$ s'identifie à la gerbe des relèvements de $P_{o}$, pour
l'homom.\ $j : \underline{H}' \to [\underline{\mathbb{Z}}]$ associé ;
$\gamma : \underline{H}' \to (\hat{\underline{\mathbb{Z}}})_{\underline{\mathbb{Q}}_{p}}$
…}

\page{78}
de passage au quotient (sauf erreur de signe, car je n'ai nulle part fait
attention aux signes).

8.4. \quad Considérons maintenant l'homomorphisme
\[
(8.4.1) \qquad \underline{G}(\bar{\underline{\mathbb{F}}}_{p})
\longrightarrow \underline{G}(\underline{\mathbb{F}}_{p}) .
\]
Il est clair qu'il se factorise par $\underline{H}$ (i.e.\ que le composé
$\underline{\mathbb{Q}}_{p}$-faisceaux constructibles sur
$\underline{\mathbb{F}}_{p} \xrightarrow{\ i\ }
\mathrm{FCriso}(\underline{\mathbb{F}}_{p}) \longrightarrow
\mathrm{FCriso}(\bar{\underline{\mathbb{F}}}_{p})$ s'envoie dans les objets
« constants » …), d'où un homomorphisme
\[
(8.4.2) \qquad \underline{G}(\bar{\underline{\mathbb{F}}}_{p})
\longrightarrow \underline{H} .
\]
Par Dieudonné-Manin, on connait la structure des objets de
$\mathrm{FCriso}(\bar{k})$ pour $\bar{k}$ algébriquement clos, on sait que
ces objets sont semi-simples et proviennent (mod isomorphisme) d'objets de
$\mathrm{FCriso}(\underline{\mathbb{F}}_{p})$. Cela implique que (8.4.1),
donc aussi (8.4.2), est un monomorphisme de gerbes. Pour montrer que c'est
une équivalence, il suffit de montrer qu'un objet simple non constant de
\struck{\ill{}} $\operatorname{Rep}(\underline{H})$ donne un objet non
constant de $\operatorname{Rep}(\underline{G}(\bar{\underline{\mathbb{F}}}_{p}))
\simeq \mathrm{FCriso}(\bar{\underline{\mathbb{F}}}_{p})$. Or les classes
d'objets simples de $\operatorname{Rep}(\underline{H})$ sont indexées par les
éléments de $\underline{\mathbb{Q}}$, et si on considère l'objet simple de
$\mathrm{FCriso}(\underline{\mathbb{F}}_{p}) \simeq
\operatorname{Rep}(\underline{G}(\underline{\mathbb{F}}^{*}_{p}))$ défini par
$(E, F_{E})$, avec $F_{E}$ ayant comme polynôme minimal $T^{r} - p^{s}$ avec
$r, s \in \underline{\mathbb{Z}}$, $r \geq 1$, $(r,s) = 1$, alors sa
restriction à $\underline{H}$ donne un objet simple \struck{\ill{}} (car il
reste simple dans $\mathrm{FCriso}(\bar{\underline{\mathbb{F}}}_{p})$) de
classe $s/r \in \underline{\mathbb{Q}}$, et celui-ci, d'après la
classification de Dieudonné-Manin, \struck{\ill{}} a une image dans
$\mathrm{Fc\,iso}(\bar{\underline{\mathbb{F}}}_{p})$ qui n'est pas constante
si $r \neq 0$ …

8.5. \quad On a donc \struck{déterminé} résolu par l'affirmative, pour
$k = \underline{\mathbb{F}}_{p}$, la question 6.2, en déterminant en même
temps (à équivalence définie à isomorphisme unique près) la gerbe
$\underline{G}(\bar{k})$ associée à une clôture algébrique $\bar{k}$ de
$k = \underline{\mathbb{F}}_{p}$ ; \struck{\ill{}} elle ne dépend pas (à
\struck{\ill{}} équivalence déterminée à isomorphisme unique près) du choix
de la clôture algébrique envisagée $\bar{k}$, et par Dieudonné-Manin est
can.\ équi-

\page{79}
valente à $\underline{G}(\gamma)$ pour tout corps alg.\ clos $\gamma$ de
car.\ $p > 0$ (de façon précise, si $\bar{\underline{\mathbb{F}}}_{p}$ est la
clôture algébrique de $\underline{\mathbb{F}}_{p}$ dans $\gamma$, l'inclusion
de celle-ci dans $\gamma$ induit une équivalence de gerbes
\[
\underline{G}(\gamma) \approx \underline{G}(\bar{\underline{\mathbb{F}}}_{p})
\add{\simeq \underline{H}} \quad ).
\]

8.6. \quad Etude de $\operatorname{Rep}(\underline{H})$. Nous avons dit que
les classes d'objets simples sont en correspondance biunivoque avec les
éléments de $\underline{\mathbb{Q}}$ ; l'élément de $\underline{\mathbb{Q}}$
correspondant à un objet isotypique est dit la \emph{pente} de cet
élément. Nous avons vu que l'objet de $\operatorname{Rep}(\underline{H})$
défini par l'objet $\underline{\mathbb{Q}}_{p}[T]/(T^{r} - p^{s})$ de
$\mathrm{FCriso}(\underline{\mathbb{F}}_{p})$, muni de la multiplication par
$T$, est simple de pente $s/r$ ($s, r \in \underline{\mathbb{Z}}$,
$r \geq 1$, $(s,r) = 1$), de sorte que nous obtenons ainsi un catalogue
complet des objets simples $E_{r,s}$. Bien entendu, $E_{r,s}$ est de rang
$r$, d'autre part son corps des endomorphismes est un corps gauche
\add{$K_{r,s}$} dont la classe dans
$\mathrm{Br}(\underline{\mathbb{Q}}_{p}) = \underline{\mathbb{Q}}/\underline{\mathbb{Z}}$
est $s/r \bmod 1$, d'après ce qu'on a vu dans 8.3.
$\operatorname{Rep}(\underline{H})$ est limite inductive des
sous-catégories pleines $\operatorname{Rep}(\underline{H}, r)$, $r \geq 1$,
où $(H,r)$ est le quotient de la gerbe $\underline{H}$ correspond au quotient
$(H,r) \approx \underline{G}_{m}$ de son lien, tel que l'homomorphisme de
Tate \struck{\ill{}} induise un homomorphisme
$(H,r) \xrightarrow{\ ?\ } \underline{G}_{m}$ s'identifiant à
$x \mapsto x^{r}$ : \add{$F$-isocristaux (sur $\bar{k}$) de pente multiple de
$1/r$.} Donc $(H,r)$ \struck{\ill{}} est la gerbe de lien
$\underline{G}_{m}$, ayant comme classe dans
$\mathrm{Br}(\underline{\mathbb{Q}}_{p})$ le générateur canonique
$1/r \bmod \underline{\mathbb{Z}}$ de
$_{r}\mathrm{Br}(\underline{\mathbb{Q}}_{p}) \simeq
\underline{\mathbb{Z}}/r\underline{\mathbb{Z}}$. C'est aussi, plus
canoniquement, « la » gerbe définie par le corps gauche $K_{r,1}$ plus haut,
\struck{mais}

\marginal{Notation $E_{s/r}$ pour $E_{r,s}$ indiquée, donc
$E_{\lambda} \otimes E_{\lambda'} \simeq
E_{\lambda\lambda'}^{n(\lambda,\lambda')}$ pour
$\lambda, \lambda' \in \underline{\mathbb{Q}}$.}
\note{Dans le membre de droite, l'indice se lit $\lambda\lambda'$ sur la
page ; la somme $\lambda + \lambda'$ que le contexte appellerait n'est pas
écrite, et on laisse comme il est.}

\add{[}Exercice amusant, que je n'ai pas fait : la structure additive de
$\operatorname{Rep}(\underline{H})$ étant entièrement décrite par la
collection des $K_{r,s}$, expliciter sa $\otimes$-structure …\ i.e.\
écrire les « relations canoniques » en termes de produits tensoriels entre
les $E_{r,s}$. Autre exercice suggéré : trouver « canoniquement » une
description des modules de Dieudonné des réduits des groupes de Lubin-Tate en
termes des $E_{r,s}$ (en particulier, les Lubin-Tate associés aux extensions
non ramifiées

\page{80}
de $\underline{\mathbb{Q}}_{p}$) qui sont donc isomorphes
\add{« géométriquement » i.e.\ dans $\operatorname{Rep}(\underline{H})$,}
(mais sans doute non canoniquement) aux $E_{r,1}$ pour $r \geq 1$ : ils
donnent des modules libres de rang $1$ canoniques sous les
$K_{r,1}$ …\add{]}

9. \quad Pour $k$ quelconque, on trouve un \struck{\ill{}}
$\otimes$-homomorphisme canonique \struck{\ill{}} défini à isomorphisme
unique près (on utilise un choix d'une clôture algébrique $\bar{k}$ de $k$,
mais ce choix est inessentiel …)
\[
\mathrm{FCriso}(k) \longrightarrow \operatorname{Rep}(\underline{H})
\]
ou encore un homomorphisme de gerbes
\[
\underline{H} \longrightarrow \underline{G}(k)
\]
(rappelons que cet homomorphisme se factorise par le noyau, \struck{\ill{}}
soit \struck{\ill{}} $\underline{H}(k)$, de
$\underline{G}(k) \to \underline{\Pi}_{p}(k)$, et que pour $k$ parfait on
devine que l'homomorphisme obtenu
$\underline{H} \to \underline{H}(k)$ est une équivalence …). Si
$\lambda \in \underline{\mathbb{Q}}$, un $F$-isocristal sur $k$ est dit
\struck{\ill{}} \emph{de pente} $\lambda$ si son image dans
$\operatorname{Rep}(\underline{H})$ est de pente $\lambda$. On voit alors
aussitôt qu'un $F$-isocristal est de pente zéro sss il est « trivial » i.e.\
satisfait aux conditions de 5.2 (ce qui justifie la terminologie qui est
introduite, un peu prématurément, à cet endroit), il est effectif sss son
image dans $\operatorname{Rep}(H)$ l'est, i.e.\ est à pente $\geq 0$.

Considérons l'homomorphisme de liens associé
\[
H \longrightarrow G(k) ,
\]
on trouve alors que si $k$ est parfait, l'homomorphisme précédent est
\emph{central}, en d'autres termes $H$ opère sur le foncteur identique
de $\mathrm{FCriso}(k) \to \operatorname{Rep}(\underline{G}(k))$, de sorte
que tout $F$-isocristal \add{$(E, F_{E})$} se décompose en ses
« \emph{composantes isopentiques} » $E_{\lambda}$,
$\lambda \in \underline{\mathbb{Q}}$. $E$ est « trivial » resp.\ effectif sss
$E_{\lambda} = 0$ pour $\lambda \neq 0$ (resp.\ $\lambda < 0$).

Notons que le \add{$F$-iso}cristal de Tate $K(1)$ est de pente $-1$. Donc
(pour $k$ parfait) tout $F$-isocristal se met de façon unique sous la forme
\[
E = \sum_{i \in \underline{\mathbb{Z}}} E_{i}(-i) ,
\]
où $E_{i}$ est à composantes isopentiques de pente $0 \leq i < 1$.
D'ailleurs,
\note{Le tapuscrit s'interrompt là, en bas de la dernière page du lot.}

\end{document}
